\begin{abstract}
\setstretch{1.08}
\setlength{\parskip}{1.5pt}
Báo cáo kỹ thuật này trình bày thiết kế, hiện thực hóa và quy trình đánh giá thực nghiệm hệ thống \textbf{NewsLens} --- giải pháp hỏi đáp tin tức dựa trên mô hình Sinh có Tăng cường Truy xuất (Retrieval-Augmented Generation -- RAG). Trên kho dữ liệu gồm 11.064 bài báo CNN và 1.152 câu hỏi bổ sung ngữ cảnh tìm kiếm từ tập dữ liệu NewsQA, bài toán đặt ra thách thức lớn từ hai nguồn sai sót độc lập: tầng truy xuất bỏ sót đoạn văn bản chứa bằng chứng, và tầng sinh câu trả lời bịa đặt thông tin (hallucination) hoặc trích dẫn nguồn không chính xác.

Hệ thống NewsLens kết hợp ứng dụng hỏi đáp (frontend React, backend FastAPI, cơ sở dữ liệu vector Chroma và chỉ mục thưa BM25S) với quy trình thực nghiệm đa giai đoạn được đăng ký trước (pre-registration). Mọi quyết định lựa chọn cấu hình đều tuân thủ nguyên tắc kiểm định bootstrap ghép cặp gom cụm theo bài báo (paired article-cluster bootstrap) kết hợp hệ thống rào chắn an toàn (safety guardrails) đa chỉ số, đảm bảo cấu hình được chọn duy trì chất lượng ổn định và cân bằng trên các tiêu chí đánh giá thay vì chỉ ưu tiên điểm số của một chỉ số đơn lẻ.

Kết quả thực nghiệm trên tập phát triển và tập kiểm định độc lập chỉ ra:
\begin{enumerate}[noitemsep, topsep=1pt, parsep=0pt]
  \item \textbf{Tầng truy xuất (Giai đoạn 1):} Mô hình BGE-M3 Sparse vượt trội các mô hình dense retriever với mức tăng $+0{,}1634$ nDCG@5 (CI95 $[+0{,}1231;\,+0{,}2080]$). Việc bổ sung bộ tái xếp hạng Cross-Encoder BGE-Reranker-Large tiếp tục cải thiện $+0{,}0659$ nDCG@5, đưa tỷ lệ tìm thấy bằng chứng Hit@5 đạt $95{,}73\%$ trên tập phát triển và $89{,}78\%$ trên 871 câu hỏi kiểm định cuối kỳ.
  \item \textbf{Tầng sinh đáp án (Giai đoạn 2):} Prompt P2 (ràng buộc loại thông tin trong một câu duy nhất) kết hợp 5 đoạn ngữ cảnh vượt qua cả 4 rào chắn an toàn, trong khi P2-depth3 bị loại do vi phạm rào chắn Faithfulness ($-0{,}0200$) và Citation Validity ($-0{,}0142$). Thử nghiệm mở rộng cho thấy prompt one-shot (P2-1S) cải thiện rõ rệt độ khớp hình thức (Exact Match tăng lên $0{,}3559$, Token F1 đạt $0{,}6267$, Answer Correctness đạt $0{,}8015$). Khảo sát các mô hình sinh làm rõ sự đánh đổi: Gemini 3.5 Flash-Lite đạt Answer Correctness cao nhất ($0{,}8164$) và độ trễ P95 thấp nhất ($0{,}89\text{ s}$), trong khi Gemini 3.1 Flash-Lite duy trì độ trung thực (Faithfulness $0{,}9603$) và chất lượng trích dẫn tốt nhất.
  \item \textbf{Bài học phương pháp luận:} Xuyên suốt ba giai đoạn thử nghiệm liên tiếp (Phase 2B, Phase 2C và Phase 3), các cấu hình đạt điểm số cao nhất ở chỉ số chính đều bị loại do vi phạm rào chắn an toàn đã đăng ký trước, khẳng định vai trò của hệ thống rào chắn đa chỉ số trong việc ngăn ngừa rủi ro suy giảm hiệu năng trên các khía cạnh khác.
  \item \textbf{Khả năng tổng quát hóa:} Cấu hình tối ưu khóa trước đạt Answer Correctness $0{,}7157$ (zero-shot) và $0{,}7557$ (one-shot) trên 284 câu hỏi của 50 bài báo held-out hoàn toàn chưa từng dùng để tinh chỉnh, chứng minh tính ổn định và độ tin cậy của hệ thống trên dữ liệu thực tế.
\end{enumerate}

\noindent\textbf{Mã nguồn (GitHub):} \url{https://github.com/ThomasdeCarpio/Text-Mining---NewsQA-RAG/tree/main} \\
\textbf{Video Demo:} \url{https://drive.google.com/file/d/1M3nf3z-kyqBO_QBA6DVe4uClQBiRqwQG/view?usp=sharing}

\noindent\textbf{Từ khóa:} Retrieval-Augmented Generation (RAG), NewsQA, Dense \& Sparse Retrieval, Cross-Encoder Reranking, One-shot Prompting, Generator Ablation, Citation Verification, Pre-registration.
\end{abstract}
